\chapter{Cleansing Techniques: Dhauti, Neti and Kapalabhati}
\label{ch:cleansing}

\emph{Which cleansing practices must the candidate know, and which one must the candidate perform?}
The syllabus asks for knowledge of dhauti and neti and for the practice of kapalabhati. The difference matters: dhauti and neti are learned only from a qualified teacher and are asked about in the exam, while kapalabhati is demonstrated. This chapter treats kapalabhati in full and the others by meaning and type. The idea of shatkarma as a whole is in the introduction to the cleansing actions (see \cref{ch:shatkarma-intro}).

\section{Kapalabhati}

\subsection*{Meaning}
\emph{Kapala} is the skull and \emph{bhati} is shining, so the name means ``skull shining''. It is a breathing practice classed with the cleansing actions, which is why it is also called kapalabhati pranayama. Its special feature is that it reverses ordinary breathing. Normally the in-breath is the active part and the out-breath is passive. In kapalabhati the out-breath is active, made by a quick pull of the abdominal muscles, and the in-breath happens by itself when the muscles relax.

\subsection*{Method}
\begin{enumerate}
\item Sit in a meditative posture such as sukhasana, padmasana or vajrasana. Close the eyes and relax the body.
\item Inhale deeply through both nostrils and let the chest expand.
\item Breathe out forcefully by contracting the abdominal muscles, then relax them so that the next in-breath comes on its own.
\item Continue as rapid breaths: active out-breath, passive in-breath. The chest and shoulders do not move.
\item The protocol gives one round as 30 rapid breaths, followed by a deep breath in and a slow breath out. Do three rounds, each followed by deep breathing.
\end{enumerate}
Beginners may start with up to three rounds of 20 breaths and build up slowly. Do not strain at any stage.

\begin{caution}
Avoid kapalabhati with a heart condition, giddiness, high blood pressure, vertigo, chronic nose bleeding, epilepsy, migraine, stroke, hernia or gastric ulcer. Check this list before every class that includes it.
\end{caution}

\subsection*{Teaching it}
Because it is quick and forceful, kapalabhati is the practice most often done too hard. Teach the abdominal pull slowly first, with attention on the sound of the out-breath, and add speed only when the pull is clean. A common fault is a heaving chest and lifting shoulders; a hand on the lower abdomen helps the student feel where the movement should come from.

\ycbfig{ycb1-fig-kapalbhati}

\section{Dhauti}

Dhauti is the cleansing of the digestive tract up to the stomach. The guidebook names three common forms:
\begin{itemize}
\item \textbf{Vaman dhauti}, with lukewarm salted water.
\item \textbf{Danda dhauti}, with a soft tube.
\item \textbf{Vastra dhauti}, with a strip of cloth.
\end{itemize}
For the exam, know the meaning and the types. Dhauti is learned only from a qualified teacher and is never taught from a book.

\ycbfig{ycb1-fig-dhauti}

\section{Neti}

Neti is the cleansing of the nasal passages. Two types are named:
\begin{itemize}
\item \textbf{Jala neti}: lukewarm salted water poured through one nostril and out of the other, using a neti pot.
\item \textbf{Sutra neti}: a soft thread or a thin rubber tube passed through the nose and out of the mouth.
\end{itemize}
Like dhauti, neti is learned only from a qualified teacher.

\ycbfig{ycb1-fig-neti}

\section{Other cleansing actions known by name}

Three further actions belong to the same family. They are outside the practice list, and a candidate should know them by name and purpose only.
\begin{itemize}
\item \textbf{Nauli}: churning and massage of the abdomen by separating and moving the central abdominal muscles. The guidebook gives four forms: madhya (middle), vama (left), dakshina (right) and nauli chalana (rotation).
\item \textbf{Trataka}: steady gazing at a small mark without blinking, traditionally until the eyes water. It is taught for the eyes and for concentration. Neither the Hatha Pradipika nor the Gheranda Samhita splits it into inner and outer forms.
\item \textbf{Basti}: cleansing of the colon; jala basti uses water and sthala basti uses air.
\end{itemize}

\ycbfig{ycb1-fig-nauli}
\ycbfig{ycb1-fig-tratak}
\ycbfig{ycb1-fig-basti}

\begin{recall}
\begin{itemize}
\item Practise kapalabhati; know dhauti and neti.
\item Kapalabhati: active out-breath by the abdomen, passive in-breath; chest and shoulders still.
\item Protocol: three rounds of 30 breaths; beginners up to three rounds of 20.
\item Dhauti types: vaman, danda, vastra. Neti types: jala, sutra.
\item Nauli forms: madhya, vama, dakshina, chalana. Trataka: steady gazing. Basti: colon.
\end{itemize}
\end{recall}

\subsection*{Worked examples}
\begin{itemize}
\item \textbf{A beginner's chest and shoulders heave with every breath.} Put a hand on the lower abdomen and practise slow pulls with a relaxed chest. Add speed only when the chest stays still.
\item \textbf{A student with high blood pressure wants to join the kapalabhati round.} Do not include them. Offer gentle deep breathing instead and suggest they ask their physician.
\item \textbf{Model answer: explain why kapalabhati is called a cleansing practice and a breathing practice.} It is listed among the cleansing actions, and it works through the breath, with a forceful out-breath and a passive in-breath.
\end{itemize}

\subsection*{Check yourself}
\begin{enumerate}
\item Which phase of the breath is active in kapalabhati?
\item Give the protocol rounds for kapalabhati, and the beginner's rounds.
\item Name four conditions in which kapalabhati is avoided.
\item Name the types of dhauti.
\item Name the two types of neti.
\item Name the forms of nauli.
\item What is trataka?
\item How are dhauti and neti learned, and what does the exam ask about them?
\end{enumerate}

\subsection*{Answers}
\begin{enumerate}
\item The out-breath, by the abdominal muscles; the in-breath is passive.
\item Three rounds of 30 breaths; beginners up to three rounds of 20.
\item Any four of: heart condition, giddiness, high blood pressure, vertigo, chronic nose bleeding, epilepsy, migraine, stroke, hernia, gastric ulcer.
\item Vaman, danda and vastra dhauti.
\item Jala neti and sutra neti.
\item Madhya, vama, dakshina and nauli chalana.
\item Steady gazing at a small mark without blinking, for the eyes and concentration.
\item Only from a qualified teacher; the exam asks for knowledge (meaning and types), not performance.
\end{enumerate}

\paragraph{Sources.} Common Yoga Protocol, Ministry of AYUSH (kapalabhati); the official guidebook for Level I (2016, older scheme), chapter on the shat kriyas; the Board syllabus for Level 1.
